% ===== PRÉAMBULE CANONIQUE — à coller VERBATIM en tête de CHAQUE .tex =====
\documentclass[11pt,a4paper]{article}
\usepackage[utf8]{inputenc}\usepackage[T1]{fontenc}\usepackage{lmodern}
\usepackage[french]{babel}
\usepackage{amsmath,amssymb,mathtools}\usepackage{icomma}
\usepackage[version=4]{mhchem}          % \ce{...} : équations & formules
\usepackage{siunitx}\sisetup{locale=FR} % \qty{}{}, \si{}, \ang{}
\usepackage{chemfig}                    % \chemfig{...} : structures développées
\usepackage{xcolor}\usepackage{colortbl}
\usepackage{tikz}\usepackage{pgfplots}\pgfplotsset{compat=1.18}
\usepackage[most]{tcolorbox}\usepackage{enumitem}\usepackage{array,booktabs}
\usepackage[a4paper,margin=2.2cm]{geometry}
\usetikzlibrary{arrows.meta,calc,angles,quotes,positioning,patterns,backgrounds,shapes.geometric}
\definecolor{primary}{RGB}{222,49,99}\definecolor{primarydark}{RGB}{150,22,55}
\definecolor{defcol}{RGB}{0,114,178}\definecolor{methcol}{RGB}{52,92,148}
\definecolor{excol}{RGB}{0,135,90}\definecolor{attncol}{RGB}{200,40,55}
\tcbset{boxbase/.style={enhanced,breakable,boxrule=0.9pt,arc=2.5pt,
  left=3mm,right=3mm,top=2.3mm,bottom=2.3mm,fonttitle=\bfseries,coltitle=white}}
% --- boîtes TUTORIEL (noms EXACTS, ne pas renommer) ---
\newtcolorbox{cle}[1][]{boxbase,colback=defcol!6,colframe=defcol,
  title={\textsc{À retenir}\ifx\relax#1\relax\relax\else\textmd{ : \itshape #1}\fi}}
\newtcolorbox{methode}[1][]{boxbase,colback=methcol!7,colframe=methcol,sharp corners,
  title={\textsc{Méthode}\ifx\relax#1\relax\relax\else\textmd{ --- \itshape #1}\fi}}
\newtcolorbox{exemplebox}[1][]{boxbase,colback=excol!5,colframe=excol,
  title={\textsc{Exemple}\ifx\relax#1\relax\relax\else\textmd{ : \itshape #1}\fi}}
\newtcolorbox{piege}[1][]{boxbase,colback=attncol!6,colframe=attncol,
  title={\textsc{Piège}\ifx\relax#1\relax\relax\else\textmd{ : \itshape #1}\fi}}
\newtcolorbox{remarque}[1][]{boxbase,colback=black!4,colframe=black!45,
  title={\textsc{Remarque}\ifx\relax#1\relax\relax\else\textmd{ : \itshape #1}\fi}}
% --- boîtes EXERCICE / CORRIGÉ (noms EXACTS, auto-comptées) ---
\newtcolorbox[auto counter]{exo}[1][]{boxbase,colback=primary!4,colframe=primary,
  title={\textsc{Exercice}~\thetcbcounter\ifx\relax#1\relax\relax\else\textmd{ --- \itshape #1}\fi}}
\newtcolorbox[auto counter]{corrige}[1][]{boxbase,colback=excol!4,colframe=excol,
  title={\textsc{Corrigé}~\thetcbcounter\ifx\relax#1\relax\relax\else\textmd{ --- \itshape #1}\fi}}
\newcommand{\R}{\mathbb{R}}\newcommand{\N}{\mathbb{N}}\newcommand{\Z}{\mathbb{Z}}\newcommand{\ds}{\displaystyle}
% =========================================================================
\begin{document}
\begin{corrige}[enthalpie d'une combustion par la loi de Hess]
\begin{enumerate}[label=\alph*)]
  \item $\Delta H^\circ_{\text{r}} = \big[\Delta H^\circ_{\text{f}}(\ce{CO2}) + 2\,\Delta H^\circ_{\text{f}}(\ce{H2O}(l))\big] - \big[\Delta H^\circ_{\text{f}}(\ce{CH4}) + 2\underbrace{\Delta H^\circ_{\text{f}}(\ce{O2})}_{=0}\big]$.
  \item $\Delta H^\circ_{\text{r}} = [(-393) + 2(-286)] - (-75) = (-393 - 572) + 75 = \qty{-890}{\kilo\joule}$ (exothermique).
\end{enumerate}
\end{corrige}

\begin{corrige}[inverser et diviser une réaction]
$\Delta H$ est proportionnel à l'équation.
\begin{enumerate}[label=\alph*)]
  \item réaction inverse : signe opposé, $\Delta H = \qty{+200}{\kilo\joule}$.
  \item coefficients divisés par $2$ : $\Delta H = \dfrac{-200}{2} = \qty{-100}{\kilo\joule}$.
\end{enumerate}
\end{corrige}

\begin{corrige}[remonter à une enthalpie de formation]
\begin{enumerate}[label=\alph*)]
  \item $\ce{H2}$ et $\ce{O2}$ sont des corps purs simples ($\Delta H^\circ_{\text{f}}=0$), donc
  \[ \Delta H^\circ_{\text{r}} = 2\,\Delta H^\circ_{\text{f}}(\ce{H2O}(l)) - 0 = 2\,\Delta H^\circ_{\text{f}}(\ce{H2O}(l)). \]
  \item $\Delta H^\circ_{\text{f}}(\ce{H2O}(l)) = \dfrac{-572}{2} = \qty{-286}{\kilo\joule\per\mol}$.
\end{enumerate}
\end{corrige}

\begin{corrige}[combiner deux réactions]
\begin{enumerate}[label=\alph*)]
  \item on additionne (1) et l'\emph{inverse} de (2) :
  \[ (1) + [-(2)] : \ce{C}(s) + \ce{O2} + \ce{CO2} \rightarrow \ce{CO2} + \ce{CO} + \tfrac12\,\ce{O2}, \]
  soit, après simplification, $\ce{C}(s) + \tfrac12\,\ce{O2}(g) \rightarrow \ce{CO}(g)$.
  \item $\Delta H = \Delta H_1 - \Delta H_2 = -393 - (-283) = \qty{-110}{\kilo\joule}$.
\end{enumerate}
\end{corrige}

\begin{corrige}[quand l'état physique change]
\begin{enumerate}[label=\alph*)]
  \item $\Delta H^\circ_{\text{r}} = \Delta H^\circ_{\text{f}}(\ce{H2O}(g)) - \Delta H^\circ_{\text{f}}(\ce{H2O}(l)) = -242 - (-286) = \qty{+44}{\kilo\joule\per\mol}$.
  \item oui : $\Delta H^\circ_{\text{r}}>0$ (endothermique), la vaporisation \emph{absorbe} de l'énergie, comme attendu pour un passage liquide $\rightarrow$ gaz.
\end{enumerate}
\end{corrige}

\end{document}
